%% kapa: exam=1 round=20 year=2009 period=1 form=A subject=economics subject_ko=경제원론 seq=27 number=67
%% source: Q-Net 감정평가사 기출문제 (cst00309) artlSeq=5024178 file=extracted/1008/5024178/2009년 제20회 감정평가사 1차시험 기출문제/1교시.hwp
%% conversion: kordoc 4.24.51 -> tools/gen1.py (auto); review=proofread-sonnet; flags=none
\begin{question}{67}
통화량 증대의 효과가 \uline{아닌} 것은?
\begin{choices}
  \item 화폐의 가치는 일반적으로 하락한다.
  \item 장기적으로 물가는 상승한다.
  \item 고전학파의 이분법에 따르면 명목소득은 증가한다.
  \item 화폐의 중립성이 성립하면 실질이자율이 하락한다.
  \item 고전학파 이론에 따르면 장기적으로 실질소득은 변화가 없다.
\end{choices}
\end{question}
\answer{4}
